\documentclass[man,12pt]{apa7}

\usepackage{graphicx}
\usepackage{setspace}

\begin{document}

\begin{titlepage}
\singlespacing

\begin{center}

% Encabezado superior
\begin{minipage}{0.7\textwidth}
\raggedright
\textbf{\large Universidad Tecnológica Nacional}\\
Facultad Regional Rosario
\end{minipage}
\begin{minipage}{0.25\textwidth}
\raggedleft

\includegraphics[width=2.2cm]{imagenes/UTN_logo.jpg}
\end{minipage}

\vspace{1.5cm}

{\Huge \textbf{Gestión Gerencial}}

\vspace{0.5cm}

{\Large Trabajo Práctico N°0}

\vspace{0.8cm}

{\LARGE \textbf{Análisis de la empresa:}}

\vspace{1cm}

\end{center}

\noindent
\textbf{Docentes:}\\
Dr. Ing. Hernán Cornejo\\
Esp. Esteban Cordiviola

\vspace{0.4cm}

\textbf{Comisión:}\\
503

\vspace{0.4cm}

\textbf{Año de cursado:}\\
2026

\vspace{0.6cm}

\textbf{Integrantes:}

\vspace{0.3cm}

\begin{tabular}{|c|l|l|}
\hline
\textbf{Legajo} & \textbf{Nombre y Apellido} & \textbf{E-mail} \\
\hline
 & Berli, Nahuel & \\
\hline
 & Jaca, Juan Pablo & \\
\hline
 & Giampietro, Gustavo & \\
\hline
51922 & Mondino, Juan Cruz & Juancm.2000@hotmail.com \\
\hline
\end{tabular}

\vfill

\end{titlepage}

\doublespacing % volvés a APA

\section{1. Análisis del grado de complejidad de la empresa}

\subsection{Matriz de Complejidad}

\begin{center}
\begin{tabular}{|l|c|c|c|}
\hline
\textbf{Dimensión} & \textbf{Alta} & \textbf{Media} & \textbf{Baja} \\
\hline
Ciclo de vida del sector &  & X &  \\
\hline
Cantidad y naturaleza de empresas del sector &  & X &  \\
\hline
Arquitectura empresarial &  & X &  \\
\hline
Naturaleza del producto o servicio & X &  &  \\
\hline
Tipo de clientes & X &  &  \\
\hline
Grado de innovación &  & X &  \\
\hline
Dependencia de insumos extranjeros &  &  & X \\
\hline
Grado de calificación RRHH & X &  &  \\
\hline
Coyuntura del sector &  & X &  \\
\hline
\end{tabular}
\end{center}

\subsection{Ciclo de vida del sector}

La digitalización de las pymes agroindustriales recién ahora está en crecimiento, lo cual da tranquilidad que no hay competidores muy grandes pero a su vez cualquier problema o error lleva a la organización que lo comete a la ruina debido a que no hubo otras que ya se hayan enfrentado a este anteriormente.

Es decir, la tranquilidad de que no haya tantos competidores se contrasta con la incertidumbre del ciclo de vida creciente del sector, por ende, definimos la complejidad de la dimensión del ciclo de vida del sector como media.

\subsection{Cantidad y naturaleza de empresas del sector}

En nuestra opinión la complejidad de esta dimensión es media.

Por más que el sector sea América del sur, si un sistema desarrollado en EEUU o Europa es mejor (ya sea por utilidad, precio, prestaciones, etc), nada impide su implementación ya que esa es la ventaja de los sistemas, que si les generan ganancia, pueden ser aplicados a lo largo del mundo.

La única posible ventaja que puede llegar a tener frente a competencias tan grandes es que este sistema tenga funcionalidades específicas del área, como por ejemplo, generación de reportes para entidades gubernamentales, cálculo de impuestos, etc.

\subsection{Arquitectura empresarial}

Definimos la arquitectura empresarial como media.

La estructura de la organización es simple y sin grandes complicaciones. Está conformada por un ingeniero en sistemas, un técnico agrónomo, un administrador de empresas, personal de integración, capacitación y soporte. Además, teniendo en cuenta su importante crecimiento, deducimos que seguramente contadores, personal de marketing y atención al cliente.

Pero todo esto sin tener un esquema organizacional definido, lo que lleva a la empresa a tener menos agilidad frente a cambios y una jerarquía no definida; a su vez, al no ser una red tan grande de diferentes puestos con nombres muy específicos o propios de la empresa, es algo más fácil el entendimiento de la “estructura”.

\subsection{Naturaleza del producto o servicio}

Esta dimensión la consideramos de complejidad alta.

Es una organización que utiliza tanto tecnologías simples (Excel para el manejo de datos) como complejas (diferentes herramientas o sensores para obtener información). Las herramientas que cada cliente puede utilizar están dadas por el plan de pago elegido, por lo que se debe tener un buen sistema de permisos para cada tipo de usuario.

Además, otorga un servicio de soporte y capacitación para ayudar a los empleados de la misma a comenzar a utilizar el sistema.

\subsection{Tipo de clientes}

Esta dimensión consideramos que tiene una complejidad alta.

Este pensamiento está fundado en que al crear una herramienta de seguimiento para este sector hay que tener en cuenta que los clientes no solo hablarán el idioma español, sino que en la zona también se habla portugués y hay que adaptar la UI de la app, herramientas o página web para diferentes idiomas. Esto se debe a que en la región se habla mayoritariamente en español o que mediante el sistema se podría informar mediante gráficos BI a inversores o sectores de gestión manejados por personal extranjero que hable al menos en inglés.

Además, el hecho de que esté dirigido a pymes o empresas que todavía no se digitalizaron y quieren dar el primer paso, genera complicaciones en capacitación en uso del sistema. Generalmente, en grandes empresas, la decisión de digitalización está dada por el área de gestión y los productores u operarios puede que no estén del todo de acuerdo y se resistan a la implementación y cambio.

\subsection{Grado de innovación}

La complejidad en el grado de innovación es media.

Las herramientas de gestión de información no son del todo complejas al hablar de tecnologías, pero al ser una organización que te introduce a la digitalización no debería tardar en incorporar compatibilidad con sensores IoT, inteligencia artificial o modelos predictivos que están muy en auge este último tiempo gracias a su gran utilidad.

\subsection{Dependencia de insumos extranjeros}

Es una empresa que brinda un producto de software y un servicio de soporte o capacitación para el mismo, por ende la dependencia de insumos extranjeros es de complejidad baja.

\subsection{Grado de calificación RRHH}

La empresa debe tener un recurso humano con una formación tal que le permita desarrollar la aplicación (sea móvil o web): analistas que elaboren todos los artefactos necesarios para que el grupo de desarrolladores entiendan cómo deben funcionar las funcionalidades específicas del sector agropecuario, desarrolladores, personal que realice test de las funcionalidades, gente de marketing, comercial o ventas; contadores, ingenieros, veterinarios, personal de soporte, personal de capacitación, consultores de las diferentes áreas que abarcan y de nuevas tecnologías, personal de RRHH, ingenieros que decidan la mejor estructura de servidores y comunicación para almacenar y transportar información, etc.

Es por ello que determinamos que la complejidad es alta.

\subsection{Coyuntura del sector}

Al ubicarse principalmente en América Latina, región la cual tiene cambios políticos constantes (que pueden traducirse en impuestos nuevos o fluctuación en el precio de la moneda, por ejemplo). Pero estos cambios llevan tiempo, no van a ser impuestos de la noche a la mañana, es por ello que determinamos que ésta dimensión tiene una complejidad media.

\subsection{Conclusión final}

Al analizar la organización como un todo, determinamos un grado de complejidad media. Ya que al interactuar todas estas dimensiones se pueden ir disminuyendo la complejidad entre sí para las que tienen una complejidad alta y el caso contrario para las dimensiones de complejidad baja.

\subsection{PESTLE}

\textbf{P - Político}
\begin{itemize}
    \item Muchos cambios en las políticas agronómicas (como puede ser el cambio en el porcentaje de retenciones), afectan a los clientes y el poder adquisitivo que tienen para adquirir nuevos servicios.
    \item Los constantes cambios en las políticas de Argentina (al igual que en el resto de América Latina) generan incertidumbre a la hora de tomar decisiones y complica la planificación de la empresa.
    \item Políticas que incentiven la exportación de servicios permiten que la empresa pueda crecer en el mercado vendiendo servicios al exterior con bajos o nulos aranceles.
\end{itemize}

\textbf{E - Economía}
\begin{itemize}
    \item Al encontrarse en América Latina, el contexto económico que lo rodea es fluctuante debido a los cambios radicales de ideologías en los gobiernos electos, lo que se ve impactado en el tipo de cambio de la moneda de cada país con respecto al dólar (moneda utilizada en el mercado internacional).
    \item En momentos de crisis o incertidumbre, los clientes priorizan los gastos operativos para tratar de seguir como estaban hasta el momento, más que invertir en la digitalización de la organización. Esto obstaculizaría la adopción del sistema.
    \item Puede que los bancos limiten el acceso a planes de financiamiento, los cuales este tipo de organizaciones necesitan para la incorporación de nuevas tecnologías.
    \item Aparición de nuevos impuestos: afectan a la organización en la parte de rentabilidad, volviendo más complicado que saquen rédito de lo que generan. No es tan común que los impuestos sean removidos.
\end{itemize}

\textbf{S - Socio-culturales}
\begin{itemize}
    \item Resistencia al cambio: los productores agrícolas están acostumbrados a no utilizar la tecnología por lo que es difícil que acepten la digitalización.
    \item Al no estar acostumbrados a la tecnología digital, es muy probable que se les dificulte el uso del sistema.
\end{itemize}

\textbf{T - Tecnología}
\begin{itemize}
    \item Implementación de la inteligencia artificial.
    \item Aparición de máquinas agrícolas que contengan sensores IoT para automatización.
    \item Problemas de conectividad en zonas rurales.
\end{itemize}

\textbf{L - Legislación}
\begin{itemize}
    \item Protección de datos debido a las leyes de seguridad de la información.
    \item Puede haber normativas que requieran trazabilidad, sanidad o control productivo de los cultivos, que obliguen organizaciones a digitalizarse lo que genera nuevos posibles clientes. Pero, también es una exigencia para AgroSys, ya que debe cumplir con dichas normativas, si no lo hace queda afuera del mercado o puede recibir sanciones y reducir la confianza de los clientes para con ella.
    \item Constante actualización de normativas de entes recaudadores (ARCA/AFIP) para las facturaciones y liquidaciones.
    \item Las regulaciones laborales dificultan la ampliación de la empresa a nivel de personal ya que incrementan altamente el costo de contratación de empleados.
\end{itemize}

\textbf{E - Ecología}
\begin{itemize}
    \item Adopción de un modelo más sustentable en la producción agrícola para el cual se debe desarrollar una herramienta específica.
    \item Tendencia a reducir el uso de agroquímicos que implica un mejor manejo del stock de la organización en ciertos productos.
\end{itemize}

\subsection{Cadena de valor}

\includegraphics{imagenes/Cadena de valor.png}

\textbf{Margen}

\textbf{Generan valor:}
\begin{itemize}
    \item Software de gestión especializado en el agro.
    \item Acompañamiento en la digitalización.
    \item Capacitación y soporte contínuo.
\end{itemize}

\textbf{Generan costos:}
\begin{itemize}
    \item Desarrollo del software.
    \item Soporte técnico.
    \item Implementación.
    \item Mantenimiento del sistema.
\end{itemize}

\subsection{Análisis de Fuerzas Competitivas}

\subsubsection{Rivalidad y competencia del mercado}

La competencia en el sector AgTech está creciendo por lo que las empresas deben seguir innovando para poder retener y conseguir nuevos clientes.

Aunque el sector está en crecimiento, muchas pymes agroindustriales se oponen al cambio por lo que las distintas empresas deben competir por aquellas pymes que ya tomaron la decisión de digitalizarse.

Los costos de ingresar al mercado son altos ya que se requiere de una inversión inicial importante y luego hay un costo en la infraestructura (servidores) si se la organización adopta un esquema SaaS. Por lo tanto, se deben captar el máximo de suscriptores posibles para afrontar esos costos.

No existe mucha diferenciación entre los software para las empresas agricultoras por lo que las empresas se deben diferenciar en la capacitación y el soporte que se le da a los clientes.

Los competidores suelen tener diferentes orígenes por lo que seguramente tengan diferentes estrategias sobre la manera de competir.

El sector tiene algunas barreras de salida. Algunas de ellas son el grado de especialización que tiene el software desarrollado para la agroindustria por lo que tienen poco valor para otros mercados y la gran la identificación de los socios con el sector por lo que continúan compitiendo incluso cuando la rentabilidad no es buena.

\subsubsection{Amenaza de nuevos productos/servicios sustitutos}

Un producto sustituto puede llegar a ser las planillas de Google Sheets + IA integrada, ya que una organización del sector agroindustrial puede preferir utilizarlas por su bajo costo y no contar con un software integral.

\subsubsection{Amenaza de nuevos competidores}

La entrada de un nuevo competidor al mercado de SaaS para la agroindustria tendrá un costo unitario mayor por cliente ya que al principio tendrá una pequeña cantidad de clientes y el costo del mantenimiento de la infraestructura será repartido entre menos clientes comparado a las organizaciones que ya tienen varios clientes.

Además, para ingresar al sector se debe hacer una inversión importante en la infraestructura junto con los sueldos de los desarrolladores de software. A su vez, las empresas que ya cuentan con el software de AgroSys fueron capacitadas en base al software que ofrece, si quisieran cambiar a otro software tendrían que capacitar a todos los empleados en la forma de utilizar el nuevo software.

Por último, AgroSys ya tiene experiencia en el sector desde su fundación, habiendo pasado por crisis y pudiendo haberlas superado. En cambio, un nuevo competidor no cuenta con esa experiencia por lo que seguramente se encontrará con problemas los cuales AgroSys ya sabe cómo solucionarlos.

\subsubsection{Poder de negociación de los clientes}

El poder de negociación de los clientes es bajo a pesar de que el software va dirigido a las pymes las cuales no tienen alta rentabilidad por lo que presionan para reducir el precio del mismo. Uno de los factores más importantes es que la calidad de los productos agrícolas de las pymes se ve muy afectada por el software ofrecido, de esta manera los clientes son mucho menos sensibles a los precios. A su vez, los clientes no pueden comparar a AgroSys con otras organizaciones del sector ya que AgroSys incorpora servicios de implementación, capacitación y soporte continuo que los demás competidores no tienen tan desarrollado. Por último, el costo de cambiar de organización proveedora del software es muy alto ya que el cliente deberá capacitar nuevamente a los empleados para que aprendan a utilizar el nuevo software.

\subsubsection{Poder de negociación de los proveedores}

El poder de negociación de los proveedores se considera alto. Existen pocas empresas que sean proveedores de infraestructura en la nube, por lo cual tienen un gran poder de negociación al poder fijar precios debido a la falta de otras opciones. También, si AgroSys quiere cambiar la plataforma de servidor, la migración implica usar una gran parte del tiempo por lo que el costo del cambio es elevado. Adicionalmente, AgroSys representa una fracción muy chica para las empresas que ofrecen servidores en la nube lo que no le da poder de negociación con ellas.

\subsection{Factores Críticos de Éxito}

\subsubsection{Adaptación del producto a las necesidades del usuario}

El sistema que la empresa vende debe ser lo suficientemente intuitivo y fácil de usar para que las pymes del sector se adapten al cambio, además debemos ofrecer una capacitación al cliente. Lo que debemos medir en este caso es el nivel de satisfacción del usuario, de esta forma la medida será los usuarios activos en el sistema sobre el total de usuarios que adquirieron el software.

\subsubsection{Innovación tecnológica continua}

El sector está en crecimiento por lo que constantemente debemos competir con las demás empresas que ofrecen servicios similares. Para mantenernos como una empresa competitiva en el sector debemos ofrecer mejoras del sistema a lo largo del tiempo, puede ser agregar nuevas funcionalidades al sistema o mejorar aquellas que ya tenemos. Para medir este FCE cada vez que se mejora el sistema debemos medir en los meses próximos si el porcentaje de nuevos clientes se mantiene de acuerdo a lo esperado, o si varía (crece o decrece).

\subsubsection{Focalización en un nicho específico}

Las pymes agroindustriales suelen tener diferentes actividades, por lo que el sistema también podría buscar definirse y consolidarse mayormente en una actividad específica. Por ejemplo, si la actividad agrícola está creciendo más que la ganadera, y vemos que hay más clientes de ese sector, podemos mejorar toda la parte agrícola del sistema, sin descuidar la parte ganadera pero focalizándonos en aquel sector. Para medir este punto, una vez que publicitamos las mejoras en este aspecto del sistema, podemos ver la tasa de nuevos clientes en un área determinada, midiendo los nuevos clientes del área X sobre el total de nuevos clientes.

\end{document}